% This file was converted to LaTeX by Writer2LaTeX ver. 1.9.9
% see http://writer2latex.sourceforge.net for more info
\documentclass{article}
\usepackage{calc,amsmath,amssymb,amsfonts}
\usepackage[LGR,T1]{fontenc}
\usepackage[greek,italian]{babel}
\usepackage[style=numeric,backend=biber]{biblatex}
\author{HP}
\date{2026-08-18}
\begin{document}
\section{L$\text{\textgreek{’}}$automobile comincia a muoversi}
Alla fine del passaggio precedente avevamo ottenuto qualcosa che, fino a quel momento, era esistito soltanto nella
nostra immaginazione: l$\text{\textgreek{’}}$automobile era comparsa sullo schermo.

Era ancora poca cosa rispetto al gioco che avevamo in mente. Non c$\text{\textgreek{’}}$era una strada, non esistevano
avversari, non avevamo ostacoli, punteggio o suoni. E soprattutto quella piccola automobile rossa rimaneva
perfettamente immobile, quasi fosse stata parcheggiata al centro del monitor in attesa che decidessimo che cosa farne.

Eppure il risultato aveva già un significato preciso. Avevamo trasformato un disegno in 63 byte, collocato quei byte in
memoria, configurato il VIC-II e ottenuto finalmente qualcosa di visibile. Per la prima volta il progetto aveva assunto
una forma.

Adesso dovevamo compiere un passo molto più importante.

Dovevamo permettere al giocatore di intervenire.

Un videogioco, infatti, non nasce nel momento in cui compare qualcosa sullo schermo. Comincia davvero a prendere vita
quando ciò che vediamo reagisce alle nostre azioni. Da quel momento non osserviamo più soltanto il risultato di un
programma: entriamo nel programma, anche se indirettamente, e cominciamo a modificarne continuamente lo stato.

La domanda per questa nuova fase è quindi semplice da formulare:

come facciamo a spostare l$\text{\textgreek{’}}$automobile verso destra e verso sinistra?

La risposta ci porterà molto più lontano di quanto potrebbe sembrare.

\subsection{Una posizione che deve poter cambiare}
Nella M01 avevamo stabilito la posizione iniziale dello sprite scrivendo direttamente le coordinate nei registri del
VIC-II:

poke 53248,172

poke 53249,120

Il primo valore determinava la posizione orizzontale, il secondo quella verticale. Per
un$\text{\textgreek{’}}$automobile immobile non avevamo bisogno di altro.

Adesso, però, il valore 172 non può più rimanere semplicemente incorporato in un$\text{\textgreek{’}}$istruzione.

Se il giocatore si sposta verso sinistra, la posizione dovrà diminuire. Se si muove verso destra, dovrà aumentare.
Qualche istante dopo potrebbe cambiare ancora.

Abbiamo quindi bisogno di conservare quella posizione in qualcosa che possa assumere valori differenti durante
l$\text{\textgreek{’}}$esecuzione del programma.

È precisamente il compito di una variabile.

La nuova inizializzazione comincia così:

240 rem posizione e limiti laterali

250 x=172:h=0:oh=0

260 xl=24:xr=320

Per il momento concentriamoci sulla prima assegnazione:

x=172

Abbiamo scelto X come nome della variabile che rappresenterà la posizione orizzontale
dell$\text{\textgreek{’}}$automobile.

Una variabile può essere immaginata come un contenitore identificato da un nome. Il programma può memorizzarvi un
valore, leggerlo quando ne ha bisogno e, soprattutto, sostituirlo con un valore differente.

All$\text{\textgreek{’}}$inizio abbiamo:

X = 172

ma durante il gioco potremmo ottenere:

X = 175

poi:

X = 178

e successivamente tornare a:

X = 175

La differenza rispetto alla M01 è sostanziale. Prima la posizione apparteneva direttamente
all$\text{\textgreek{’}}$istruzione POKE; adesso appartiene allo stato del nostro programma.

Il VIC-II continuerà naturalmente ad avere bisogno della coordinata nei propri registri, ma sarà il BASIC a conservarne
il valore corrente e a decidere quando modificarlo.

Nella stessa riga inizializziamo altre due variabili:

h=0:oh=0

Il carattere : permette di scrivere più istruzioni BASIC sulla stessa riga. Il Commodore le eseguirà
nell$\text{\textgreek{’}}$ordine in cui le incontra, come se fossero state distribuite su linee separate.

Il significato di H e OH diventerà chiaro tra poco, quando scopriremo che la coordinata orizzontale di uno sprite
possiede una particolarità che la nostra variabile X, almeno per il momento, riesce a nasconderci molto bene.

La riga successiva introduce invece altre due variabili:

xl=24:xr=320

I nomi sono stati scelti per ricordarne facilmente il significato: possiamo leggere XL come X Left, il limite sinistro,
e XR come X Right, il limite destro.

Stiamo quindi dichiarando fin dall$\text{\textgreek{’}}$inizio che l$\text{\textgreek{’}}$automobile potrà muoversi
nell$\text{\textgreek{’}}$intervallo compreso tra 24 e 320.

Avremmo potuto inserire questi due numeri direttamente più avanti, dentro le istruzioni che controllano il movimento. Il
programma avrebbe funzionato allo stesso modo. Tenerli invece insieme alla configurazione iniziale ci offre un
vantaggio importante: stiamo separando la regola dai valori che la configurano.

Oggi i nostri limiti sono 24 e 320. Quando costruiremo davvero la strada potremmo scoprire che la carreggiata deve
essere più stretta. In quel caso sarà sufficiente cambiare:

xl=24:xr=320

senza cercare all$\text{\textgreek{’}}$interno del ciclo principale tutte le istruzioni nelle quali quei numeri sono
stati utilizzati.

È una scelta semplice, ma introduce già un$\text{\textgreek{’}}$abitudine utile: quando un valore rappresenta una
caratteristica del gioco destinata potenzialmente a cambiare, conservarlo in una variabile rende il programma più
leggibile e più facile da modificare.

Vale la pena ricordare anche una particolarità del BASIC V2: nei nomi delle variabili sono significativi soltanto i
primi due caratteri. XL e XR sono quindi perfettamente distinti, mentre nomi più lunghi che iniziassero con le stesse
due lettere verrebbero considerati equivalenti. È uno dei limiti del linguaggio con cui dovremo imparare a convivere
scegliendo nomi brevi, ma sufficientemente riconoscibili.

Definite posizione e limiti, possiamo inizializzare realmente lo sprite:

270 poke 53248,x

280 poke 53249,120

290 m=peek(53264)and254

300 poke 53264,m

La coordinata verticale rimane ancora fissa a 120. In questa fase vogliamo affrontare un problema alla volta e ci
occuperemo esclusivamente del movimento laterale.

\subsection{Un programma che non deve più finire}
Finora il nostro listato era costituito soprattutto da una fase di preparazione. Puliva lo schermo, caricava i dati
dello sprite, configurava il VIC-II e mostrava l$\text{\textgreek{’}}$automobile.

Un programma interattivo non può però limitarsi a fare tutto questo una sola volta.

Deve continuare a domandarsi se il giocatore stia facendo qualcosa, decidere se la posizione debba cambiare, aggiornare
eventualmente lo sprite e poi ricominciare.

Nasce così quella che nel listato abbiamo chiamato esplicitamente:

520 rem ciclo principale

Da questo punto il programma entra in una sequenza destinata a ripetersi continuamente:

530 wait 53266,128

540 wait 53266,128,128

550 k=peek(197)

560 if k=10 then x=x-s:goto 590

570 if k=18 then x=x+s:goto 590

580 goto 530

Il GOTO non è un concetto completamente nuovo. Molti anni prima Rossi Brunori lo aveva utilizzato per mostrarci come
costringere il Commodore a ripetere continuamente un semplice PRINT {\textquotedbl}CIAO{\textquotedbl}.

Allora avevamo scoperto che il computer non si stanca.

Adesso quella stessa proprietà assume un significato completamente diverso.

Non vogliamo ripetere continuamente la stessa azione. Ogni volta che il ciclo ricomincia dobbiamo osservare ciò che sta
accadendo e decidere quale comportamento adottare.

Il nostro programma non deve quindi soltanto ripetersi.

Deve reagire.

\subsection{Prima di muoverci, decidiamo quando farlo}
Le prime due istruzioni del ciclo possono sembrare piuttosto enigmatiche:

530 wait 53266,128

540 wait 53266,128,128

È la prima volta che incontriamo WAIT ed è anche il primo momento in cui il nostro gioco comincia a preoccuparsi non
soltanto di che cosa debba fare, ma anche di quando farlo.

Per comprenderne il motivo dobbiamo affacciarci, senza ancora entrarvi troppo in profondità, sul funzionamento del
sistema video.

Il VIC-II non costruisce l$\text{\textgreek{’}}$intera immagine che vediamo sul monitor nello stesso istante. Procede
progressivamente, linea dopo linea. Possiamo immaginare una sorta di scansione che attraversa continuamente lo schermo
e che prende il nome di raster.

L$\text{\textgreek{’}}$indirizzo 53266 ci permette di osservare una parte della posizione corrente di questa scansione.

Non abbiamo ancora bisogno di studiare dettagliatamente il raster. Lo faremo quando il gioco ci presenterà problemi per
i quali sarà realmente necessario comprenderlo meglio. Per ora ci interessa soltanto una proprietà: il suo movimento ci
offre un riferimento che si ripete regolarmente.

Ed è proprio quello che utilizziamo con WAIT.

La forma più semplice dell$\text{\textgreek{’}}$istruzione è:

wait indirizzo,maschera

WAIT sospende l$\text{\textgreek{’}}$esecuzione del BASIC finché i bit che ci interessano nella locazione indicata non
raggiungono una determinata condizione.

Nel nostro caso utilizziamo come maschera il numero 128.

In binario:

128 = 10000000

Stiamo quindi osservando il bit 7.

Con:

wait 53266,128

aspettiamo che quel bit diventi 1.

Subito dopo:

wait 53266,128,128

utilizziamo il terzo parametro di WAIT per attendere la condizione opposta, cioè il ritorno dello stesso bit a zero.

La coppia di istruzioni crea per noi un punto di sincronizzazione che si presenta una volta per ciascun ciclo completo
del video. Possiamo quindi utilizzarlo come una sorta di metronomo.

Il ciclo non aggiornerà l$\text{\textgreek{’}}$automobile semplicemente ogni volta che il BASIC riesce a completare
tutte le istruzioni. Prima aspetterà il riferimento video che abbiamo scelto.

È ancora una tecnica molto elementare. Non stiamo utilizzando interrupt, non stiamo controllando con precisione ogni
linea raster e non pretendiamo ancora di sincronizzare il gioco come farebbe un programma interamente scritto in
linguaggio macchina.

Abbiamo semplicemente iniziato a far sì che il nostro programma ascolti anche il ritmo della macchina sulla quale viene
eseguito.

\subsection{PEEK: questa volta vogliamo leggere}
Dopo aver atteso il momento previsto dal ciclo troviamo:

550 k=peek(197)

Qui incontriamo un$\text{\textgreek{’}}$istruzione che svolge, in un certo senso, l$\text{\textgreek{’}}$operazione
opposta rispetto a uno strumento che abbiamo già utilizzato molte volte.

Con POKE scriviamo:

poke indirizzo,valore

Con PEEK, invece, leggiamo il contenuto di una locazione di memoria:

peek(indirizzo)

La nostra istruzione:

k=peek(197)

può quindi essere letta così:

leggi il valore contenuto nella locazione 197 e conservalo nella variabile K.

La locazione 197 viene utilizzata dal sistema per registrare il codice associato alla tastiera. Il nostro programma può
quindi osservare quel valore per capire se il giocatore sta premendo uno dei comandi che abbiamo deciso di utilizzare.

Il percorso concettuale è questo:

giocatore

\ \ \ \ ↓

tastiera

\ \ \ \ ↓

locazione 197

\ \ \ \ ↓

PEEK(197)

\ \ \ \ ↓

variabile K

Il Commodore non sa che stiamo cercando di guidare un$\text{\textgreek{’}}$automobile.

Per lui esistono soltanto numeri.

Saremo noi ad attribuire a quei numeri un significato.

\subsection{IF: il programma comincia a decidere}
Subito dopo la lettura troviamo:

560 if k=10 then x=x-s:goto 590

570 if k=18 then x=x+s:goto 590

IF introduce uno dei concetti fondamentali della programmazione: la possibilità di eseguire
un$\text{\textgreek{’}}$istruzione soltanto quando una determinata condizione è vera.

La struttura generale è:

if condizione then istruzione

e può essere letta quasi come una frase naturale:

se questa condizione è vera, allora fai questo.

Consideriamo:

if k=10 then x=x-s:goto 590

Il programma verifica innanzitutto:

K = 10 ?

Se la risposta è negativa, passa semplicemente alla riga successiva.

Se invece la condizione è vera, esegue:

x=x-s

e poi:

goto 590

È interessante osservare come il simbolo = assuma due significati leggermente diversi all$\text{\textgreek{’}}$interno
della stessa istruzione.

In:

k=10

stiamo verificando se K sia uguale a 10.

In:

x=x-s

stiamo invece assegnando a X un nuovo valore.

Se, per esempio:

X = 172

S = 3

dopo:

x=x-s

avremo:

X = 169

La riga successiva effettua il movimento nella direzione opposta:

if k=18 then x=x+s:goto 590

Questa volta aggiungiamo S.

Se nessuna delle due condizioni è verificata, arriviamo a:

580 goto 530

e torniamo direttamente all$\text{\textgreek{’}}$inizio del ciclo.

Questa scelta è importante. Se il giocatore non ha richiesto alcuno spostamento, non abbiamo motivo di ricalcolare e
riscrivere la posizione dello sprite.

Non è cambiato nulla.

Il Commodore può quindi tornare ad aspettare il riferimento video successivo.

È una delle prime piccole attenzioni alle prestazioni del nostro programma: non eseguire un lavoro quando sappiamo già
che non produrrà alcun cambiamento.

\subsection{Quanto deve spostarsi?}
Poco prima del ciclo principale abbiamo inizializzato:

490 rem 3 pixel per fotogramma

500 s=3

La variabile S rappresenta l$\text{\textgreek{’}}$entità dello spostamento applicato alla coordinata X ogni volta che il
giocatore richiede un movimento.

Questo semplice valore nasce in realtà da una scelta che merita di essere spiegata.

Lo spostamento ideale, dal punto di vista della continuità visiva, sarebbe di un solo pixel alla volta. In questo modo
ogni nuova posizione dello sprite sarebbe adiacente alla precedente e il piccolo salto percepibile durante il movimento
verrebbe ridotto al minimo.

Avremmo, per esempio:

172

173

174

175

176

177

Con il BASIC V2, tuttavia, questa scelta rende lo spostamento laterale dell$\text{\textgreek{’}}$automobile troppo lento
per il tipo di gioco che stiamo costruendo.

Aumentare l$\text{\textgreek{’}}$incremento rende invece la macchina più reattiva.

Con:

s=3

potremmo ottenere:

172

175

178

181

184

L$\text{\textgreek{’}}$automobile percorre quindi una distanza maggiore nello stesso numero di aggiornamenti, ma tra una
posizione e la successiva vengono saltate due coordinate.

Questo introduce inevitabilmente un piccolo compromesso visivo.

La macchina reagisce più velocemente ai comandi, ma il movimento non possiede la continuità spaziale che otterremmo
procedendo pixel per pixel.

Per il momento accettiamo quindi questo compromesso e utilizziamo uno spostamento di tre pixel.

Il lieve micro-scattare che ne deriva rimane consapevolmente come un limite della soluzione corrente. Più avanti, quando
il gioco ci porterà ad affrontare l$\text{\textgreek{’}}$Assembly, potremo tornare sul problema disponendo di un
controllo molto più preciso sui tempi di aggiornamento e sulla gestione del movimento.

È importante sottolineare che non stiamo scegliendo tre perché sia un numero magicamente corretto.

Stiamo scegliendo il valore che, nelle nostre prove, rappresenta il miglior equilibrio tra due esigenze contrastanti:
fluidità e reattività.

Anche questo significa progettare.

Non sempre esiste una soluzione migliore in senso assoluto. Spesso esiste la soluzione più adatta agli strumenti e ai
vincoli che abbiamo in quel momento.

\subsection{Due limiti che adesso hanno un nome}
Una volta modificata la posizione dobbiamo assicurarci che l$\text{\textgreek{’}}$automobile rimanga
nell$\text{\textgreek{’}}$intervallo che abbiamo deciso di utilizzare.

Nella parte iniziale del programma avevamo scritto:

260 xl=24:xr=320

Possiamo quindi esprimere il controllo senza inserire nuovamente i due valori numerici:

590 if x{\textless}xl then x=xl

600 if x{\textgreater}xr then x=xr

La prima istruzione significa:

se X è diventata più piccola del limite sinistro XL, riportala esattamente a XL.

La seconda compie l$\text{\textgreek{’}}$operazione opposta:

se X supera il limite destro XR, riportala a XR.

Il nostro intervallo rimane quindi:

XL {\textless}= X {\textless}= XR

e, con i valori attuali:

24 {\textless}= X {\textless}= 320

Dal punto di vista del comportamento non è cambiato nulla rispetto all$\text{\textgreek{’}}$uso diretto dei numeri 24 e
320.

Dal punto di vista della struttura del programma, invece, la differenza è importante.

Il ciclo principale non ha più bisogno di conoscere quali siano i limiti. Deve soltanto sapere che esistono un limite
sinistro XL e uno destro XR.

Se domani decidessimo che il movimento deve avvenire fra 50 e 290, ci basterebbe modificare una sola riga:

xl=50:xr=290

e le istruzioni:

if x{\textless}xl then x=xl

if x{\textgreater}xr then x=xr

continuerebbero a funzionare senza alcuna variazione.

È una distinzione che incontreremo sempre più spesso: una parte del programma contiene la logica,
un$\text{\textgreek{’}}$altra contiene i valori sui quali quella logica deve lavorare.

Oggi la differenza sembra piccola.

Quando il programma crescerà, diventerà sempre più preziosa.

\subsection{Il problema del numero 320}
Proprio il limite destro ci mette però davanti a una difficoltà interessante.

Abbiamo stabilito:

xr=320

e quindi la nostra variabile X può assumere valori superiori a 255.

Ma il registro:

53248

nel quale il VIC-II conserva la coordinata orizzontale dello sprite 0 è formato da un solo byte.

Un byte contiene otto bit.

Con otto bit possiamo rappresentare 256 combinazioni:

0 ... 255

Come possiamo allora collocare lo sprite a X=256?

O a 280?

O al nostro limite 320?

La variabile BASIC non ha alcun problema a contenere questi valori. La difficoltà nasce soltanto quando dobbiamo
trasferirli al VIC-II.

La coordinata orizzontale degli sprite viene infatti rappresentata attraverso nove bit.

Gli otto bit meno significativi dello sprite 0 sono contenuti nel registro 53248.

Il nono bit si trova invece in un altro registro:

53264

Questo secondo registro è condiviso tra tutti gli otto sprite: ciascuno dei suoi bit costituisce il nono bit della
coordinata X di uno sprite diverso.

Per lo sprite 0 viene utilizzato il bit 0.

Possiamo quindi immaginare la coordinata in questo modo:

\ \ \ \ \ \ \ \ \ \ \ \ coordinata X


\bigskip

bit 8 \ \ \ \ bit 7 bit 6 bit 5 bit 4 bit 3 bit 2 bit 1 bit 0

\ \ ↓ \ \ \ \ \ \ \ \ \ \ \ \ \ \ \ \ \ \ \ \ \ \ \ \ \ \ \ ↓

53264 \ \ \ \ \ \ \ \ \ \ \ \ \ \ \ \ \ \ \ \ \ \ \ 53248

bit 0

La nostra variabile X contiene il numero completo.

Adesso dobbiamo trasformarlo nella rappresentazione che il VIC-II si aspetta.

\subsection{H e L: separiamo la coordinata}
La conversione avviene attraverso:

610 if x{\textless}256 then l=x:h=0

620 if x{\textgreater}255 then l=x-256:h=1

Abbiamo chiamato L la parte che finirà nel normale registro X e H il bit aggiuntivo.

Partiamo da un esempio semplice:

X = 172

Poiché 172 è inferiore a 256 viene eseguita:

if x{\textless}256 then l=x:h=0

ottenendo:

L = 172

H = 0

Non abbiamo bisogno del nono bit.

Proviamo invece con:

X = 300

Questa volta viene eseguita la seconda istruzione:

if x{\textgreater}255 then l=x-256:h=1

Otteniamo:

L = 300 - 256

L = 44


\bigskip

H = 1

La posizione 300 viene quindi rappresentata attraverso due valori:

L = 44

H = 1

Il significato rimane:

256 + 44 = 300

Gli otto bit contenuti in L rappresentano 44, mentre H fornisce il bit aggiuntivo che vale 256.

A questo punto possiamo aggiornare immediatamente il registro principale:

630 poke 53248,l

La parte bassa della coordinata cambia praticamente ad ogni spostamento.

Per aggiornare H, invece, dobbiamo fare un po$\text{\textgreek{’}}$ più di attenzione.

\subsection{Un registro appartiene a otto sprite}
Potremmo essere tentati di scrivere:

poke 53264,h

Per il nostro sprite sembrerebbe sufficiente.

Ma 53264 non contiene soltanto il suo nono bit.

Contiene quelli di tutti gli sprite:

bit \ \ \ \ \ \ \ 7 6 5 4 3 2 1 0

sprite \ \ \ \ 7 6 5 4 3 2 1 0

Scrivendo direttamente zero cancelleremmo anche i noni bit degli sprite da 1 a 7.

Scrivendo uno manterremmo acceso soltanto quello dello sprite 0 e spegneremmo comunque tutti gli altri.

Adesso abbiamo ancora una sola automobile.

Ma sappiamo già che non rimarrà sola per molto.

Conviene quindi imparare fin da ora a cambiare esclusivamente il bit che ci interessa.

Per farlo utilizziamo:

650 m=peek(53264)and254

Questa riga introduce un nuovo modo di ragionare.

Non stiamo più trattando il registro come un semplice numero.

Cominciamo a manipolarne direttamente i bit.

\subsection{AND e la nostra prima maschera}
Il numero 254, scritto in binario, diventa:

11111110

L$\text{\textgreek{’}}$ultimo bit è zero.

Gli altri sette sono uno.

L$\text{\textgreek{’}}$operatore logico AND confronta i bit corrispondenti di due valori. Il risultato sarà 1 soltanto
quando entrambi i bit confrontati valgono 1.

Immaginiamo che 53264 contenga:

10110101

Eseguiamo:

10110101

AND

11111110

{}-{}-{}-{}-{}-{}-{}-{}-

10110100

Che cosa è cambiato?

Soltanto il bit 0.

Gli altri sette hanno mantenuto esattamente il valore precedente.

Possiamo quindi rappresentare la maschera in forma più generale:

xxxxxxx?

AND

11111110

{}-{}-{}-{}-{}-{}-{}-{}-

xxxxxxx0

Le x indicano bit che vogliamo preservare.

L$\text{\textgreek{’}}$operazione:

m=peek(53264)and254

legge quindi il registro, forza a zero soltanto il bit dello sprite 0 e conserva il risultato nella variabile M.

A questo punto possiamo inserire il nuovo valore di H:

660 poke 53264,m+h

Se:

H = 0

il bit rimane spento.

Se:

H = 1

viene acceso.

Gli altri sette bit restano invece invariati.

Abbiamo così imparato una delle tecniche che diventeranno sempre più importanti man mano che ci avvicineremo
all$\text{\textgreek{’}}$hardware del Commodore: modificare soltanto i bit che ci appartengono, preservando gli altri.

\subsection{Ma H non cambia quasi mai}
Possiamo ancora evitare un$\text{\textgreek{’}}$operazione inutile.

Immaginiamo che l$\text{\textgreek{’}}$automobile si muova:

100 → 103 → 106 → 109 → 112

Il valore di H rimane sempre:

0

Se invece siamo nella zona:

280 → 283 → 286 → 289

H rimane:

1

Il nono bit cambia soltanto quando attraversiamo il confine fra 255 e 256.

Non avrebbe quindi molto senso leggere 53264, applicare la maschera e riscriverlo ad ogni singolo movimento.

Per questo all$\text{\textgreek{’}}$inizio avevamo creato:

oh=0

Possiamo leggere OH come old H: il precedente valore di H.

Dopo aver calcolato il nuovo valore troviamo infatti:

640 if h=oh then 680

Se H e OH sono uguali significa che il nono bit non è cambiato.

Possiamo quindi saltare direttamente alla riga 680.

Soltanto quando i due valori sono diversi eseguiamo:

650 m=peek(53264)and254

660 poke 53264,m+h

670 oh=h

Dopo aver aggiornato il registro memorizziamo anche il nuovo stato:

oh=h

In questo modo, al movimento successivo, potremo confrontarlo nuovamente.

È una piccola ottimizzazione, ma nasce da un ragionamento molto semplice:

se qualcosa non è cambiato, non abbiamo necessariamente bisogno di riscriverlo.

Questa attenzione sarà particolarmente importante nel BASIC V2. L$\text{\textgreek{’}}$interprete non è veloce e il
nostro gioco crescerà. Evitare operazioni inutili può quindi aiutarci a conservare tempo per tutto ciò che aggiungeremo
in seguito.

Non significa che dobbiamo ossessionarci con l$\text{\textgreek{’}}$ottimizzazione di ogni istruzione fin
dall$\text{\textgreek{’}}$inizio. Significa soltanto imparare a riconoscere i casi nei quali possiamo ottenere lo
stesso risultato svolgendo meno lavoro.

\subsection{Il ciclo si chiude}
L$\text{\textgreek{’}}$ultima istruzione della parte operativa è:

680 goto 530

Torniamo alle due WAIT.

E ricominciamo.

Il nostro primo ciclo di gioco può essere riassunto così:

\ \ \ \ \ \ \ attendi il riferimento video

\ \ \ \ \ \ \ \ \ \ \ \ \ \ \ \ \ \ \ \ ↓

\ \ \ \ \ \ \ \ \ \ \ \ leggi la tastiera

\ \ \ \ \ \ \ \ \ \ \ \ \ \ \ \ \ \ \ \ ↓

\ \ \ \ \ \ \ \ \ è richiesto movimento?

\ \ \ \ \ \ \ \ \ \ \ \ \ $\searrow $ \ \ \ \ \ \ \ \ \ \ \ \ $\searrow $

\ \ \ \ \ \ \ \ \ \ \ no \ \ \ \ \ \ \ \ \ \ \ \ \ \ sì

\ \ \ \ \ \ \ \ \ \ \ ↓ \ \ \ \ \ \ \ \ \ \ \ \ \ \ \ \ ↓

\ \ \ \ \ nuovo ciclo \ \ \ \ \ \ \ modifica X

\ \ \ \ \ \ \ \ \ \ \ \ \ \ \ \ \ \ \ \ \ \ \ \ \ \ \ \ ↓

\ \ \ \ \ \ \ \ \ \ \ \ \ \ \ \ \ \ \ \ controlla XL/XR

\ \ \ \ \ \ \ \ \ \ \ \ \ \ \ \ \ \ \ \ \ \ \ \ \ \ \ \ ↓

\ \ \ \ \ \ \ \ \ \ \ \ \ \ \ \ \ \ \ \ \ calcola L e H

\ \ \ \ \ \ \ \ \ \ \ \ \ \ \ \ \ \ \ \ \ \ \ \ \ \ \ \ ↓

\ \ \ \ \ \ \ \ \ \ \ \ \ \ \ \ \ \ aggiorna il VIC-II

\ \ \ \ \ \ \ \ \ \ \ \ \ \ \ \ \ \ \ \ \ \ \ \ \ \ \ \ ↓

\ \ \ \ \ \ \ \ \ \ \ \ \ \ \ \ \ \ \ \ \ \ \ nuovo ciclo

Il programma è ancora minuscolo rispetto a ciò che diventerà, ma la struttura fondamentale comincia già a essere
riconoscibile.

All$\text{\textgreek{’}}$inizio esiste una fase di preparazione, eseguita una sola volta.

Poi il programma entra in un ciclo che continua a osservare l$\text{\textgreek{’}}$esterno, modificare lo stato del
gioco e trasferire il risultato all$\text{\textgreek{’}}$hardware.

Quando arriveranno avversari, proiettili, collisioni e strada, questo ciclo inevitabilmente crescerà.

Ma il principio rimarrà sostanzialmente lo stesso.

\subsection{Il programma completo}
Possiamo adesso osservare la M02 nella sua interezza.

Gran parte dell$\text{\textgreek{’}}$inizializzazione deriva direttamente dalla M01 e non abbiamo bisogno di studiarla
nuovamente: sappiamo già perché lo sprite occupa 63 byte, perché li memorizziamo a partire da 832, perché il puntatore
contiene 13 e come vengono configurati modalità multicolor e colori.

Questa volta dobbiamo soprattutto osservare ciò che cambia: la posizione diventa uno stato modificabile, i limiti
laterali acquistano un nome, nasce un ciclo principale, leggiamo l$\text{\textgreek{’}}$input e affrontiamo per la
prima volta la coordinata X a nove bit.

10 rem ===============================

20 rem m02 - movimento del giocatore

30 rem coordinata x a 9 bit

40 rem ===============================

50 rem

60 rem pulisce lo schermo

70 print chr\$(147)

80 rem bordo nero

90 poke 53280,0

100 rem sfondo grigio chiaro

110 poke 53281,15

120 rem

130 rem carica i 63 byte dello sprite

140 rem a partire da \$0340 = 832

150 for i=0 to 62

160 read a

170 poke 832+i,a

180 next i

190 rem

200 rem blocco sprite: 832/64 = 13

210 rem puntatore sprite 0 = 2040

220 poke 2040,13

230 rem

240 rem posizione e limiti laterali

250 x=172:h=0:oh=0

260 xl=24:xr=320

270 poke 53248,x

280 poke 53249,120

290 m=peek(53264)and254

300 poke 53264,m

310 rem

320 rem nessun ingrandimento

330 poke 53271,0

340 poke 53277,0

350 rem

360 rem sprite 0 in multicolor

370 poke 53276,1

380 rem

390 rem multicolor 1 = bianco

400 poke 53285,1

410 rem multicolor 2 = nero

420 poke 53286,0

430 rem colore sprite 0 = rosso

440 poke 53287,2

450 rem

460 rem abilita lo sprite 0

470 poke 53269,1

480 rem

490 rem 3 pixel per fotogramma

500 s=3

510 rem

520 rem ciclo principale

530 wait 53266,128

540 wait 53266,128,128

550 k=peek(197)

560 if k=10 then x=x-s:goto 590

570 if k=18 then x=x+s:goto 590

580 goto 530

590 if x{\textless}xl then x=xl

600 if x{\textgreater}xr then x=xr

610 if x{\textless}256 then l=x:h=0

620 if x{\textgreater}255 then l=x-256:h=1

630 poke 53248,l

640 if h=oh then 680

650 m=peek(53264)and254

660 poke 53264,m+h

670 oh=h

680 goto 530

690 rem

700 rem ===============================

710 rem dati grafici dello sprite

720 rem 21 righe x 3 byte = 63 byte

730 rem ===============================

740 data \ \ 0, \ \ 0, \ \ 0, \ \ 0, \ \ 0, \ \ 0

750 data \ \ 3, 255, 192, \ \ 0, 170, \ \ 0

760 data \ \ 0, \ 40, \ \ 0, \ \ 1, 170, \ 64

770 data \ \ 3, 150, 192, \ \ 3, 150, 192

780 data \ \ 3, 150, 192, \ \ 0, 150, \ \ 0

790 data \ \ 0, 150, \ \ 0, \ \ 0, 170, \ \ 0

800 data \ \ 0, 190, \ \ 0, \ \ 5, 186, \ 80

810 data \ 15, 190, 240, \ 15, 170, 240

820 data \ 15, 255, 240, \ \ 2, 170, 128

830 data \ \ 2, 170, 128, \ \ 0, \ \ 0, \ \ 0

840 data \ \ 0, \ \ 0, \ \ 0

Se confrontiamo ciò che vediamo sul monitor con tutto ciò che il programma deve fare per produrlo, la sproporzione è
quasi sorprendente.

Agli occhi del giocatore accade pochissimo: premiamo un comando e l$\text{\textgreek{’}}$automobile si sposta.

Dietro quel gesto, però, il BASIC attende il proprio riferimento temporale, legge una locazione di memoria, interpreta
un codice, valuta una condizione, modifica una variabile, confronta quella variabile con due limiti, scompone una
coordinata in due parti, verifica se un singolo bit debba essere modificato e infine comunica la nuova posizione al
VIC-II.

Poi ricomincia.

Il giocatore non vede K, X, XL, XR, L, H, OH oppure M. Non vede la maschera 254, non percepisce le due istruzioni WAIT e
non sa che una parte della coordinata si trova in un registro diverso dalle altre.

Non deve saperlo.

Per lui l$\text{\textgreek{’}}$automobile deve semplicemente muoversi.

Per noi, invece, proprio ciò che rimane invisibile costituisce la parte interessante.

Da ragazzo, davanti ai videogiochi, continuavo a chiedermi come facesse una macchina a sapere che avevo premuto un
comando e a far reagire immediatamente qualcosa sullo schermo.

La domanda era sempre la stessa.

Come fa?

Adesso, almeno per questo piccolissimo pezzo del problema, possiamo cominciare a dare una risposta.

Il Commodore non sa che stiamo guidando un$\text{\textgreek{’}}$automobile. Non conosce il significato di destra e
sinistra e non può immaginare che quei valori numerici rappresentino i bordi entro i quali vogliamo muoverci.

Legge numeri.

Verifica condizioni.

Modifica bit.

Scrive altri numeri nei registri del VIC-II.

Siamo noi ad aver attribuito a tutte quelle operazioni un significato.

Ed è forse questo il cambiamento più importante rispetto alla M01. La nostra automobile non è più semplicemente un
disegno che siamo riusciti a far comparire sul monitor.

Adesso, finalmente, possiamo guidarla.


\bigskip
\end{document}
